% Code of Conduct.
%
% NOT derived from a .docx -- there was none. The set covered security
% (InfoSec) and harassment (POSH) and nothing else about how to behave.
% Conflicts of interest and outside work are the sections that matter most
% here, because an unpaid intern is free to work elsewhere and Paridyum has
% to say what that means rather than leave it unstated.
%
% WORDING IS A DRAFT FOR REVIEW BY QUALIFIED COUNSEL.
\documentclass{paridyum}
\begin{document}

\paridyumletterhead
\paridyumdateref{\dkDocDate}{\dkDocRef}
\paridyumtitle{Code of Conduct}

\textit{This code applies to everyone working at or with Paridyum, including
interns, for the duration of their engagement.}

\section*{1. What This Covers}
This code sets out how Paridyum expects people to conduct themselves. It sits
alongside, and does not replace, the Information Security \& Acceptable Use
Policy, the Policy on Prevention of Sexual Harassment, and the Non-Disclosure
Agreement.

\section*{2. Treating People Decently}
Everyone at Paridyum is entitled to be treated with respect. The following are
not acceptable, whether in person, in writing, or in any online space used for
work:
\begin{itemize}[leftmargin=*, noitemsep]
    \item bullying, intimidation, or persistent unwelcome behaviour;
    \item discrimination or harassment on the basis of sex, gender, caste,
          religion, disability, sexual orientation, or any other personal
          characteristic;
    \item abusive, demeaning, or threatening language; and
    \item deliberately excluding a colleague from work they should be part of.
\end{itemize}
Sexual harassment is dealt with separately and more specifically by the Policy
on Prevention of Sexual Harassment, which sets out how to complain and to whom.

\section*{3. Honesty About Work}
The Intern agrees to:
\begin{itemize}[leftmargin=*, noitemsep]
    \item represent their own progress accurately, including when work is
          behind or has gone wrong;
    \item say when they do not know something, rather than guess in a way
          that others will rely on;
    \item attribute work honestly, including work produced with the
          assistance of AI tools; and
    \item raise a problem early rather than let it become someone else's
          surprise.
\end{itemize}
Reporting a mistake promptly is treated as doing the job properly, not as a
failure.

\section*{4. Availability and Communication}
The Intern is expected to work approximately the hours recorded in their
Internship Agreement, at times broadly agreed with their supervisor, and to
tell their supervisor in advance when they will be unavailable. Paridyum does
not expect a response outside agreed working hours.

\section*{5. Outside Work and Conflicts of Interest}
\textbf{The Intern is free to hold other work, study, or an internship
elsewhere.} This internship is unpaid and part-time, and Paridyum makes no
claim on the Intern's time beyond what the Internship Agreement records.

The Intern must tell Paridyum if any of the following applies, so that work can
be assigned around it:
\begin{itemize}[leftmargin=*, noitemsep]
    \item they work for, or are engaged by, a business that competes directly
          with a Paridyum product they are working on;
    \item they have a financial interest in a supplier, customer, or
          competitor of Paridyum;
    \item a close family member is in a position to be affected by a decision
          the Intern takes for Paridyum; or
    \item they are subject to an obligation to another party that would
          conflict with the Intellectual Property Assignment Agreement.
\end{itemize}
Disclosing a conflict is not a problem. Not disclosing one is.

\section*{6. Company Property and Resources}
Paridyum accounts, systems, and equipment are provided for internship work.
Incidental personal use is acceptable; using them for another business, or in a
way that puts Paridyum's information at risk, is not. The Information Security
\& Acceptable Use Policy sets the detail.

\section*{7. Speaking About Paridyum}
The Intern may say publicly that they intern at Paridyum. They may not speak on
Paridyum's behalf, commit Paridyum to anything, or disclose confidential
information or unreleased work. Where something is unclear, ask before posting.

\section*{8. Raising a Concern}
Any concern about conduct may be raised under the Grievance Policy, which
explains where to take it and what happens next. A concern raised in good faith
will not be held against the person raising it.

\section*{9. Breach}
A serious or repeated breach of this code may lead to removal of system access
or termination of the internship, as set out in the Internship Agreement.

\section*{10. Acknowledgement}
By signing below, the Intern confirms that they have read and understood this
Code of Conduct and agree to follow it.

\vspace{1.5em}
\paridyumofficers{\dkSigCeoName}{\dkSigCeoDesignation}{\dkSigCtoName}{\dkSigCtoDesignation}

\vspace{1.5em}
\paridyumcountersign{Intern}{\dkPIdentityLegalName}

\end{document}
